\section{Rôles et modèle RBAC}

\subsection{Exercice 3.1 -- Construire la matrice RBAC du service RH}
\textbf{Objectif.} Mettre en œuvre la matrice RBAC du service RH en créant les rôles métier, en leur
attribuant les privilèges minimaux et en les affectant aux utilisateurs.

\begin{center}
\small
\begin{tabular}{|l|c|c|c|c|c|}
\hline
\textbf{Rôle / privilège} & \code{CREATE SESSION} & \code{rh.departements} & \code{rh.v_annuaire} & \code{rh.employes} (cols) & \code{rh.employes} (salaire) \\
\hline
\code{R_CONNEXION} & \textbf{Oui} & -- & -- & -- & -- \\
\hline
\code{R_RH_LECTEUR} & via \code{R_CONNEXION} & \textbf{SELECT} & \textbf{SELECT} & -- & -- \\
\hline
\code{R_RH_GESTION} & via \code{R_CONNEXION} & via \code{LECTEUR} & via \code{LECTEUR} & \textbf{SELECT, INSERT, UPD}(\code{poste},\code{dept_id}) & -- \\
\hline
\code{R_RH_PAIE} & via \code{R_CONNEXION} & -- & -- & \textbf{SELECT} & \textbf{UPD}(\code{salaire}) \\
\hline
\end{tabular}
\end{center}

\shots{3_1}

\textbf{Observations et réponses.}
\begin{enumerate}
  \item \textbf{LINA} (\code{R\_RH_LECTEUR}) : accède à \code{rh.departements} (4) et \code{rh.v_annuaire}
        (8), mais \code{SELECT * FROM rh.employes} renvoie \code{ORA-00942}.
  \item \textbf{OMAR} (\code{R\_RH_GESTION}) : accède à \code{rh.employes} (8), modifie \code{poste} avec
        succès (\code{1 row updated}), mais la modification de \code{salaire} renvoie \code{ORA-41900: missing
        UPDATE privilege}. Ses droits directs sur la table et la vue ont tous été révoqués au préalable
        (\code{DBA_TAB_PRIVS} et \code{DBA_COL_PRIVS} vides pour OMAR).
  \item \textbf{SARA} (\code{R\_RH_PAIE}) : lit \code{rh.employes}, modifie \code{salaire} avec succès
        (\code{1 row updated}), mais la modification de \code{poste} renvoie \code{ORA-41900} et l'accès à
        \code{rh.departements} renvoie \code{ORA-00942}.
\end{enumerate}

\subsection{Exercice 3.2 -- Lire la hiérarchie dans le dictionnaire}
\textbf{Objectif.} Interroger le dictionnaire pour reconstituer l'arbre d'héritage des rôles et l'ensemble
des privilèges effectifs d'un utilisateur.

\shots{3_2}

\textbf{Observations et réponses.}
\begin{enumerate}
  \item \code{DBA_ROLE_PRIVS WHERE grantee = 'OMAR'} ne renvoie que \code{R_RH_GESTION}. \code{R_RH_LECTEUR}
        et \code{R_CONNEXION} n'y apparaissent pas directement car ils sont accordés au rôle \code{R_RH_GESTION}
        et non directement à l'utilisateur OMAR.
  \item La requête hiérarchique avec \code{START WITH grantee = 'OMAR' CONNECT BY PRIOR granted_role = grantee}
        déroule l'arbre complet : \code{R_RH_GESTION} $\to$ \code{R_CONNEXION} et \code{R_RH_LECTEUR} $\to$
        \code{R_CONNEXION}.
  \item La jointure avec \code{DBA_TAB_PRIVS} / \code{DBA_COL_PRIVS} et l'interrogation de
        \code{ROLE_TAB_PRIVS} fournissent la liste exhaustive des 6 privilèges objet dont dispose OMAR
        via la hiérarchie des rôles.
\end{enumerate}

\subsection{Exercice 3.3 -- Effet immédiat ou à la reconnexion ?}
\textbf{Objectif.} Analyser le comportement des sessions ouvertes lors de la modification des privilèges
d'un rôle ou de l'attribution d'un nouveau rôle.

\shots{3_3_a}
\shots{3_3_b}
\shots{3_3_c}

\textbf{Observations et réponses.}
\begin{enumerate}
  \item \textbf{Étape 1} : quand l'administrateur accorde \code{SELECT ON rh.employes} au rôle
        \code{R_RH_LECTEUR}, la session de LINA déjà ouverte voit immédiatement la table (\code{8}) sans
        se reconnecter. Les privilèges objet portés par un rôle actif sont réévalués dynamiquement à l'exécution.
  \item \textbf{Étape 2} : après révocation du droit sur \code{R_RH_LECTEUR} et attribution du rôle
        \code{R_TEST} à LINA, le \code{SELECT} dans la même session échoue (\code{ORA-00942}). Le nouveau rôle
        accordé n'est pas actif dans une session déjà ouverte. Après \code{SET ROLE ALL}, le rôle s'active
        et la requête réussit (\code{8}). Après reconnexion, la requête réussit aussi.
  \item \textbf{Conclusion} : lors du départ d'un collaborateur, révoquer ses rôles ne suffit pas ; ses
        sessions ouvertes conservent leurs rôles actifs jusqu'à leur fermeture. Il faut obligatoirement
        \textbf{tuer les sessions ouvertes} avec \code{ALTER SYSTEM KILL SESSION} pour rendre la révocation
        immédiatement effective.
\end{enumerate}

\subsection{Exercice 3.4 -- Rôle protégé et rôles par défaut}
\textbf{Objectif.} Protéger le rôle de la paie par mot de passe et gérer les rôles par défaut.

\shots{3_4}

\textbf{Observations et réponses.}
\begin{enumerate}
  \item Après \code{ALTER ROLE r_rh_paie IDENTIFIED BY "Paie#2026"}, la tentative de connexion de Sara
        renvoie \code{ORA-01045: Login denied}. Le rôle protégé par mot de passe n'est plus activé
        automatiquement à la connexion ; comme Sara ne possédait \code{CREATE SESSION} que via
        \code{R_RH_PAIE}, elle ne peut plus ouvrir de session.
  \item Pour corriger, l'administrateur accorde \code{R_CONNEXION} directement à Sara et définit
        \code{ALTER USER sara DEFAULT ROLE r_connexion}. Sara se connecte alors avec le seul rôle
        \code{R_CONNEXION} actif ; l'accès à \code{rh.employes} renvoie \code{ORA-00942}.
  \item Sara tente \code{SET ROLE r_rh_paie IDENTIFIED BY "mauvais"} $\to$ \code{ORA-01979: missing or
        invalid password for role "R_RH_PAIE"}. Avec le bon mot de passe (\code{Paie#2026}), le rôle
        s'active (\code{Role set.}) et Sara peut à nouveau lire les employés et modifier leurs salaires.
\end{enumerate}

\ecart{Dans l'énoncé 19c, le rôle par défaut restait activé sans mot de passe à la connexion. En 26ai, un
rôle protégé par mot de passe n'est pas activé automatiquement s'il est le seul vecteur de connexion, ce qui
a provoqué \code{ORA-01045} jusqu'à l'octroi direct de \code{R_CONNEXION}.}

\subsection{Exercice 3.5 -- Le piège des rôles dans le code}
\textbf{Objectif.} Comprendre pourquoi les unités PL/SQL nommées (droits du définisseur) et les vues ignorent
les privilèges reçus par rôles.

\shots{3_5}

\textbf{Observations et réponses.}
\begin{enumerate}
  \item Bien qu'OMAR puisse exécuter \code{SELECT * FROM rh.employes} en SQL direct, la compilation de la
        fonction \code{nb_employes} échoue avec \code{PLS- / ORA-00942: table or view "RH"."EMPLOYES" does
        not exist}. De même, la création de la vue \code{v_effectif} renvoie \code{ORA-41904: Missing object
        privilege on "RH"."EMPLOYES" to create a view}.
  \item \textbf{Raison} : en PL/SQL à droits du définisseur (défaut) et pour la création de vues, le SGBD
        désactive les rôles lors de la compilation pour garantir la stabilité et la sécurité du code.
        Seuls les privilèges accordés \textbf{directement} à l'utilisateur sont pris en compte.
  \item L'option \code{AUTHID CURRENT_USER} a aussi échoué à la compilation en 26ai : le contrôle du SQL
        statique à la compilation exige que l'objet soit accessible au définisseur.
  \item \textbf{Correction la plus restreinte} : \code{GRANT SELECT ON rh.employes TO omar}. Ce privilège
        direct permet à la fonction de compiler sans erreur (\code{Function created.}) et à la vue d'être
        créée (\code{View created.}). C'est une exception assumée au principe RBAC pur.
\end{enumerate}

\subsection{Exercice 3.6 -- Que contiennent CONNECT et RESOURCE ?}
\textbf{Objectif.} Examiner le contenu des rôles système prédéfinis \code{CONNECT} et \code{RESOURCE}.

\shots{3_6}

\textbf{Observations et réponses.}
\begin{enumerate}
  \item \code{ROLE_SYS_PRIVS} montre : \code{CONNECT} contient \code{CREATE SESSION} et \code{SET CONTAINER}
        (2 privilèges). \code{RESOURCE} contient 15 privilèges de création (\code{CREATE TABLE}, \code{CREATE
        PROCEDURE}, \code{CREATE SEQUENCE}, \code{CREATE TRIGGER}, \code{CREATE TYPE}, etc.), mais
        \textbf{pas de quota} sur aucun tablespace.
  \item L'utilisateur \code{TEST_RES} doté de \code{CONNECT, RESOURCE} crée sa table \code{t_res}, mais
        l'insertion d'une ligne échoue avec \code{ORA-01950: insufficient quota on tablespace USERS}. Il
        lui manque un quota explicite (\code{UNLIMITED TABLESPACE} ayant été retiré de \code{RESOURCE}).
  \item Ces rôles historiques ne doivent pas servir de base à une politique RBAC d'entreprise : ils sont
        trop larges, figés, accordent des droits inutiles (ex: \code{CREATE CLUSTER}, \code{CREATE OPERATOR})
        et varient selon la version de la base sans correspondre aux métiers de l'organisation.
\end{enumerate}

\subsection{Exercice 3.7 -- Analyse des privilèges (DBMS\_PRIVILEGE\_CAPTURE)}
\textbf{Objectif.} Capturer les privilèges réellement utilisés par un utilisateur avec le package système.

\shots{3_7}

\textbf{Observations et réponses.}
\begin{enumerate}
  \item La capture \code{cap_omar} (\code{type => G_CONTEXT}, \code{condition => 'SESSION_USER = OMAR'})
        s'est créée (\code{CREATE_CAPTURE}), activée (\code{ENABLE_CAPTURE}) et désactivée
        (\code{DISABLE_CAPTURE}) sans erreur.
  \item Après \code{GENERATE_RESULT}, la consultation des vues \code{DBA_USED_OBJPRIVS} et
        \code{DBA_UNUSED_OBJPRIVS} est fonctionnelle (colonnes \code{CAPTURE}, \code{USERNAME},
        \code{USED_ROLE}, \code{OBJ_PRIV}, \code{OBJECT_NAME}).
  \item Cet outil permet en production d'analyser l'activité réelle d'un compte sur une période donnée pour
        retirer les privilèges inutilisés et appliquer rigoureusement le principe du moindre privilège.
\end{enumerate}
